\chapter{2008 Nobel Prize in Economics}
\textbf{Laureates: }Paul Krugman (USA)

\section*{Reason for Award}
New trade theory and new economic geography

\section{What They Did}
Paul Krugman developed new trade theory, explaining intra-industry trade, economies of scale, and imperfect competition in international markets. His models showed that even identical countries trade because consumers value variety and firms benefit from scale economies. This challenged the classical comparative advantage framework, which predicted trade only between dissimilar economies. Krugman also pioneered new economic geography, demonstrating how transportation costs and increasing returns generate spatial concentration of economic activity. His core-periphery model showed how small differences in initial conditions can lead to unequal regional development, with industries clustering in certain areas while others remain specialized in agriculture. These insights explained urbanization patterns, industrial districts, and global trade flows.

\section{Impact on Economic Thinking}
Krugman's work transformed international economics from a discipline focused on comparative advantage to one incorporating imperfect competition, product differentiation, and scale economies. His trade models provided theoretical foundations for understanding the rapid growth of intra-industry trade among developed nations. In economic geography, his spatial models challenged the neoclassical assumption of uniform distribution, showing how market forces can generate persistent regional inequalities. His work influenced trade policy debates, demonstrating that strategic trade policies might be beneficial under certain conditions. The new economic geography also informed regional development policy, explaining why infrastructure investments can have disproportionate effects on local economies. Krugman's accessible writing style brought these ideas to broader audiences, making him one of the most influential economists of his generation.

\section{Modern Perspectives Today}
Krugman's trade and geography models remain central to international economics and regional science. Modern trade theory incorporates firm heterogeneity, building on Krugman's insights about variety and scale. The global value chain revolution has validated his predictions about intermediate goods trade and production fragmentation. In economic geography, his spatial models inform urban planning, transportation policy, and cluster development strategies. The rise of China and other emerging economies has renewed interest in his trade theories, particularly regarding economies of scale and learning-by-doing. Digital platforms and e-commerce have created new spatial dynamics that extend his core-periphery framework. Climate change and sustainability concerns also connect to his work on agglomeration economies and resource allocation. Krugman's intellectual legacy endures in how economists understand globalization, urbanization, and regional development.
